\begin{table*}[t]
\centering
\caption{Other NMS-free detectors on all 5,000 val2017 images, 640, batch 1
(default: mean of two runs; locked: one run). ``Naive'' is float input with sequential execution and
stream synchronization; ``full'' is uint8 input, two prefetch workers, overlap
and event synchronization. MHz is the mean sampled GPU clock under the
default governor.}
\label{tab:generality}
\small
\setlength{\tabcolsep}{2.6pt}
\begin{tabular}{llcc rrrr rr}
\toprule
& & & & \multicolumn{4}{c}{Default} & \multicolumn{2}{c}{Locked} \\
\cmidrule(lr){5-8} \cmidrule(lr){9-10}
Model & Pipeline & AP & AP$_S$ & img/s & mJ & W & MHz & img/s & mJ \\
\midrule
YOLO26s & Ultralytics & 0.4769 & 0.3020 & 25.1 & 263.5 & 6.62 & 306 & 63.0 & 194.4 \\
 & Ours, naive & 0.4770 & 0.3019 & 24.5 & 271.6 & 6.66 & 306 & 52.2 & 221.8 \\
 & Ours, full & 0.4769 & 0.3019 & 140.1 & 131.8 & 18.47 & 1019 & 140.4 & 133.3 \\ \midrule
YOLOv10n & Ultralytics & 0.3836 & 0.1901 & 30.8 & 198.4 & 6.12 & 306 & 73.0 & 144.5 \\
 & Ours, naive & 0.3837 & 0.1900 & 33.5 & 189.4 & 6.34 & 306 & 60.0 & 167.4 \\
 & Ours, full & 0.3839 & 0.1901 & 212.4 & 76.9 & 16.33 & 1018 & 215.2 & 77.8 \\
\bottomrule
\end{tabular}
\end{table*}


\Cref{tab:generality} repeats the three main configurations for
YOLO26s (10.0\,M parameters) and YOLOv10n. In every case the three
pipelines agree in AP to within $4\times10^{-4}$, and the naive
pipelines run the GPU at 306\,MHz under default governors. YOLOv10n
behaves like YOLO26n: 30.8 images/s with Ultralytics and 212.4 with the
full pipeline ($6.9\times$), at 76.9 instead of 198.4\,mJ per image.
YOLO26s has $3.9\times$ the parameters of YOLO26n and a $1.6\times$
longer engine time (6.78 versus 4.22\,ms back-to-back). Its full
pipeline is limited by the GPU rather than the CPU. It reaches
140.1 images/s ($5.6\times$ Ultralytics), 95\% of the 146.9
inferences/s that \texttt{trtexec} achieves back-to-back, and locking
clocks does not change it (140.4). Under the default governors, the full
YOLO26s pipeline outruns the Ultralytics YOLO26n predictor with locked
clocks (140.1 vs.\ 73.9 images/s) at 0.074 higher AP and 7.5\,mJ less
per image (131.8 vs.\ 139.3). This comparison changes the pipeline and
the model at once, so it does not show that the larger model is the
better choice; YOLO26n in the same full pipeline is faster still. It
shows that the larger model costs less than a naive pipeline, even one
running with locked clocks. A deployment decision that compares model
sizes in a naive pipeline therefore mostly measures the pipeline.
